\section{Alcance y Limitaciones}
\label{chap:alcance_limitaciones}

\subsection{Alcances}
\label{sec:alcances}

El producto de esta investigación es un \textbf{módulo de simulación} para APIs nativas de Capacitor, implementado como \textbf{bridge universal extensible} con patrón \tcode{Proxy}, validado en el núcleo \tcode{@capacitor/geolocation}, \tcode{@capacitor/camera} y \tcode{@capacitor/preferences}. El módulo se integra en \textbf{al menos un proyecto Ionic de referencia} (Angular o React) y permite ejecutar pruebas unitarias en Node.js mediante Jest y/o Jasmine, con perfiles configurables de éxito, error, permisos y datos límite. Junto al simulador se entrega un \textbf{procedimiento de cobertura} para elevar líneas y ramas del código que consume esas APIs, y una \textbf{capa auxiliar de orquestación basada en MCP} que automatiza el ciclo de medición (ejecución de suite, lectura de reportes y consultas estáticas opcionales). La validación documenta una \textbf{línea base experimental} y los resultados obtenidos tras la intervención en el mismo repositorio.

\subsection{Limitaciones}
\label{sec:limitaciones}

El simulador \textbf{no emula el kernel del sistema operativo} ni el \textit{firmware} del dispositivo; reproduce únicamente el \textbf{contrato TypeScript/JavaScript} que el desarrollador invoca (\tcode{getCurrentPosition}, \tcode{getPhoto}, \tcode{set}, etc.) con respuestas coherentes en Node. Por tanto, \textbf{no detecta fallos de hardware real} (sensor GPS defectuoso, cámara ocupada en el teléfono, cuotas de almacenamiento del SO) ni sustituye pruebas en dispositivo físico o emulador nativo.

El alcance tecnológico se concentra en tres plugins del ecosistema Capacitor; otros \tcode{@capacitor/*} quedan fuera del núcleo inicial o dependen del mecanismo de \textit{fallback} genérico. Factores externos pueden condicionar los resultados: acceso limitado a hardware móvil diverso, diferencias entre versiones de sistemas operativos, cambios de versión en Ionic/Capacitor o en el runner de pruebas, y disponibilidad de servicios cloud o tokens para escenarios MCP acotados. Si el proyecto de referencia tiene tooling muy personalizado, puede requerir adaptaciones no contempladas en la primera versión del artefacto.
